% !TeX program = lualatex
% !TeX encoding = utf8
% !TeX spellcheck = uk_UA
% !BIB program = biber

\documentclass[%
biblatex,
%marginversioninfo
]{ConspectBook}
\UseTblrLibrary{booktabs}
\usepackage{listings}
\setlength{\parindent}{0pt}


%% ----------------------------------------------------------------
\title{Позначення в курсі класичної електродинаміки}
\author{}
\date{}
%% ----------------------------------------------------------------

\begin{document}
\maketitle
\newpage

\chapter{Базові поняття та рівняння}

\section*{Вектори}

Вектори позначаються \textbf{жирним шрифтом} (пряме жирне накреслення):
\[
  \vect{F},\quad \vect{V},\quad \vect{j},\quad \vect{r},\quad \vect{n},\quad
  \Efield,\quad \Bfield.
\]


\section*{Тензори}

Тензори позначаються нижніми індексами: $T_{kj}$, $\delta_{kj}$, $\epsilon_{ijk}$.

\section*{Диференціальні оператори}

\begin{center}
\begin{tblr}{
colspec={QX},
row{1} = {c,m, font=\bfseries},
}
  \toprule
  Позначення & Оператор \\
  \midrule
  $\Div\vect{A}$ &  дивергенція \\
  $\Rot\vect{A}$ &  ротор (curl) \\
  $\dfrac{\partial f}{\partial t}$ &  частинна похідна за часом\\
  $\partial_i$ &  частинна похідна за координатою $x_i$  \\
  \bottomrule
\end{tblr}
\end{center}


\section*{Поля та сила}

  $\Efield$ ---  напруженість електричного поля.\\
  $\Bfield$ ---  індукція магнітного поля.\\
  $\vect{F}$ ---  сила (сила Лоренца).\\
  $\vect{f}$ ---  об'ємна густина сили\\[4pt]


\section*{Заряди та струми}


  $q$ ---  електричний заряд.\\
  $\rho(t,\vect{r})$ ---  об'ємна густина заряду.\\
  $\sigma(\xi,\eta)$ ---  поверхнева густина заряду.\\
  $n$ ---  концентрація носіїв заряду.\\
  $I_S$ ---  сила струму через поверхню $S$.\\
  $\vect{j}(t,\vect{r})$ ---  вектор об'ємної густини струму.\\
  $\vect{i}(\xi,\eta)$ ---  вектор поверхневої густини струму.\\
  $f_k(t,\vect{r},\vect{v})$ ---  функція розподілу $k$-го сорту зарядів за швидкостями.\\


\section*{Швидкості та кінематика}


  $\vect{V}$ ---  швидкість пробного заряду.\\
  $\vect{v}$ ---  швидкість носія заряду.\\


\section*{Геометрія та координати}


  $\vect{r}=\{x,y,z\}$ ---  радіус-вектор точки ---  вектор.\\
  $d^3\vect{r} \equiv dV$ ---  елемент об'єму ---  $d^3\vect{r} = dx\,dy\,dz$.\\
  $d\vect{S} = \vect{n}\,dS$ ---  орієнтований елемент поверхні ---  $\vect{n}$ --- одинична нормаль.\\
  $d\vect{r}$ (або $d\vect{l}$) ---  елемент довжини контура ---  вектор вздовж контура.\\
  $\vect{n}$ ---  одинична нормаль до поверхні ---  напрямок вказується в кожній задачі.\\
  $\vect{\tau}$ ---  одиничний тангенціальний вектор ---  вздовж лінії на поверхні.\\
  $\Omega$ ---  область інтегрування в $\mathbb{R}^3$ ---.\\
  $\partial\Omega$ ---  межа області $\Omega$ ---  замкнена поверхня.\\
  $S$ ---  (відкрита) поверхня інтегрування ---.\\
  $\partial S$ ---  межа поверхні $S$ ---  замкнений контур.\\


\section*{Енергія, імпульс та момент імпульсу поля}

  $w = \dfrac{\Efield^2+\Bfield^2}{8\pi}$ ---  густина енергії електромагнітного поля ---  скаляр.\\
  $W$ ---  повна енергія поля в об'ємі $\Omega$ ---  $W = \int_\Omega w\,d^3\vect{r}$.\\
  $\vect{\Pi} = \dfrac{c}{4\pi}[\Efield\times\Bfield]$ ---  вектор Пойнтінга (вектор) ---  густини потоку енергії.\\
  $\vect{\pi} = \dfrac{\vect{\Pi}}{c^2}$ ---  густина імпульсу поля ---  вектор.\\
  $L_i = \epsilon_{ijk}x_j\pi_k$ ---  густина моменту імпульсу поля ---  компонентний запис.\\
  $T_{kj}$ ---  максвеллівський тензор натягу ---
  $T_{kj}=\dfrac{1}{4\pi}\!\left(E_kE_j+B_kB_j\right)-\dfrac{\delta_{kj}}{8\pi}(\Efield^2+\Bfield^2)$.\\



\end{document}
